\section{Conclusions}
%结论部分，国赛少，美赛很常见
%论文中心思想的重申，研究结果（主要观点）的归纳，某些启发性的解释或考虑
%把上一个section模型评价优缺点分析放在这部分也可以
In this study, the problem of responders planning room searches after an emergency occurs has been geometrically abstracted. We then conducted a deeper graph-theoretic modeling and used two algorithms to solve this model. After simulating three different scenarios, we found complementary properties of the two algorithms. More importantly, regardless of which algorithm was used, both significantly improved computational efficiency compared to the previous approach of merely enumerating through traversal. We found that the optimized paths could balance risk, room priority, and overall efficiency, essentially improving the original strategy. Finally, we considered how different situations affect the strategy, providing dynamic parameters and planning that further altered the paths and decisions. 

At the same time, we see this as a field with great potential for future development. With more computational input, the model might be able to conduct dynamic analysis not only for types of emergencies but also considering additional factors such as information limitations and crowd evacuation. Technology can also be increasingly integrated into our model, with the participation of robots and AI potentially leading to transformative changes in its execution.
